\documentclass[12pt,spanish,letterpaper,oneside]{article}

\def\documenttitle {Rese\~na cr\'itica: concepto y fundamentaci\'on de los derechos humanos}
\def\documentsubtitle {Dignidad, reconocimiento y garantias}
\def\documentsubject {Actividad 2 - Antecedentes de los derechos
humanos}
\def\documentauthor {Martin Jonathan de la Cruz}
\def\coursename {Antecedentes de los derechos humanos}
\def\coursecode {005}
\def\universityname {Universidad Abierta y a Distancia de Mexico}
\def\universityfaculty {Licenciatura en Derecho}
\def\universitydepartment {Antecedentes de los derechos humanos}
\def\universitydepartmentimage {departamentos/UnADM}
\def\universitydepartmentimagecfg {height=1.57cm}
\def\universitylocation {Roma Norte, Ciudad de Mexico}
\def\authortable {
  \begin{tabular}{ll}
    \textbf{Alumno:} & Martin Jonathan de la Cruz \\
    Matricula: & ES2611202040 \\
    Figura docente: & Daniel Alexis Lozano Ortega \\
    Semestre/Bloque: & 2 / 2 \\
    Actividad: & 2 \\
    Semana: & 2 \\
    \multicolumn{2}{l}{Fecha de realizacion: \today} \\
    \multicolumn{2}{l}{\universitylocation}
  \end{tabular}
}

\let\pdfpagewidth\pagewidth
\let\pdfpageheight\pageheight
\input{template}
\usepackage{fontspec}
\setmainfont{arial.ttf}[Path=C:/Windows/Fonts/,BoldFont=arialbd.ttf,ItalicFont=ariali.ttf,BoldItalicFont=arialbi.ttf]
\setsansfont{arial.ttf}[Path=C:/Windows/Fonts/,BoldFont=arialbd.ttf,ItalicFont=ariali.ttf,BoldItalicFont=arialbi.ttf]
\renewcommand{\familydefault}{\sfdefault}
\usepackage{array}
\usepackage{booktabs}
\usepackage{longtable}
\usepackage{calc}
\usepackage{tikz}
\usetikzlibrary{arrows.meta,positioning,calc,babel}
\setcitestyle{authoryear,open={(},close={)}}
\providecommand{\tightlist}{\setlength{\itemsep}{0pt}\setlength{\parskip}{0pt}}
\providecommand{\pandocbounded}[1]{#1}
\providecommand{\passthrough}[1]{#1}
\setlength{\emergencystretch}{3em}
\usepackage[most]{tcolorbox}
\newtcolorbox{resenabox}{enhanced,breakable,colback=white,
  colframe=green!30!black,boxrule=0.6pt,leftrule=2pt,arc=2pt,
  left=6pt,right=6pt,top=5pt,bottom=5pt,before skip=6pt,after skip=6pt}

\def\coverwatermarkenabled {true}
\def\coverwatermarkimage {img/departamentos/UnADM.pdf}
\def\coverwatermarkopacity {0.16}
\def\coverwatermarkwidth {0.72\paperwidth}
\def\coverwatermarkxshift {0cm}
\def\coverwatermarkyshift {-0.35cm}
\newcommand{\insertcoverwatermark}{%
  \ifthenelse{\equal{\coverwatermarkenabled}{true}}{%
    \AddToShipoutPictureBG*{%
      \begin{tikzpicture}[remember picture,overlay]
        \node[opacity=\coverwatermarkopacity,inner sep=0pt,
          xshift=\coverwatermarkxshift,yshift=\coverwatermarkyshift]
          at (current page.center) {%
            \includegraphics[width=\coverwatermarkwidth]{\coverwatermarkimage}%
          };
      \end{tikzpicture}%
    }%
  }{}%
}

\begin{document}
\insertcoverwatermark
\templatePortrait
\templatePagecfg
\onehalfspacing
\templateFinalcfg
\newgeometry{margin=2.5cm}
\fontsize{12}{14.5}\selectfont
\onehalfspacing

\section{Introduccion}

Fundamentar los derechos humanos exige distinguir por que merecen
proteccion, donde se reconocen y como se garantizan. La leccion de
Canal UNED permite valorar esa relacion desde tres perspectivas;
su contraste con Mexico requiere doctrina y normas, no solo una
enumeracion de derechos \citep{unedVideo2024,cuadernilloADH2026}.

\section{Dignidad, normas y proteccion efectiva}
La dignidad orienta la proteccion de la persona; el reconocimiento
positivo establece obligaciones juridicas y las garantias permiten
reclamarlas. El recorrido historico comprende avances y retrocesos,
no una igualdad automaticamente realizada
\citep[p.~11]{alvarezIcaza2009ADH}. La Unidad 1 incorpora garantismo
y otras corrientes que complementan, sin sustituir, el analisis del
video \citep{cuadernilloADH2026,ibarraCurso2022ADH}.

\subsection{Concepto y fundamentacion de los derechos humanos}
Se realiza una resena critica del video \emph{Tema 11. Los Derechos
Humanos Concepto y Fundamentacion}, en el que se presentan el concepto
de derechos humanos, sus principales caracteristicas y las perspectivas
que explican su fundamento. La resena examina estas ideas y valora su
importancia para comprender la proteccion juridica de la persona
\citep{unedVideo2024}.

\begin{resenabox}
\noindent\colorbox{black!5}{\begin{minipage}{\dimexpr\linewidth-2\fboxsep\relax}
\textbf{Ficha de la obra}\par
\textbf{Titulo:} \emph{Tema 11. Los Derechos Humanos Concepto y
Fundamentacion (1er Curso, 1er Cuatrimestre del GRADO EN DERECHO)}.\par
\textbf{Institucion:} UNED. \textbf{Participante:} Aliuska Duardo Sanchez.\par
\textbf{Publicacion:} 23 octubre 2024. \textbf{Formato:} leccion de
Canal UNED. \textbf{Duracion:} 17 minutos y 42 segundos.
\end{minipage}}\par\medskip
\textbf{Dignidad sin concesiones: del fundamento a la garantia}\par

La leccion de Aliuska Duardo Sanchez ofrece una entrada al problema
de los derechos humanos que considero especialmente valiosa: antes de
preguntar cuales son, obliga a pensar por que corresponden a todas las
personas. Su recorrido enlaza historia, concepto y fundamentacion.
No se limita a enumerar derechos, sino que propone distinguir aquello
que justifica su existencia de las normas que permiten reclamarlos.
Esa diferencia organiza el contenido y da sentido a la comparacion
entre las tres perspectivas que presenta.

El punto de partida historico es la distancia entre privilegios y
derechos. Las concesiones medievales beneficiaban a personas o grupos
determinados; por eso no equivalen a derechos humanos universales.
La profesora explica despues la generalizacion, la universalizacion
y la internacionalizacion, vinculada con tratados y garantias que
trascienden al Estado. Encuentro util esta secuencia porque muestra
que la proteccion no surgio completa ni se extendio sin conflictos.
Sin embargo, distinguir etapas no debe hacernos confundir la amplitud
de una declaracion con el acceso efectivo a los derechos que proclama.

La aproximacion conceptual presentada a partir de Perez Luno relaciona
dignidad, libertad e igualdad con facultades e instituciones reconocidas
por el derecho. Lo importante es esa relacion: la dignidad aporta una
razon para proteger a la persona, pero la proteccion tambien necesita
instituciones. La exposicion de universalidad, inalienabilidad,
irrenunciabilidad e imprescriptibilidad refuerza la idea de que la
titularidad no depende de un favor del gobernante. A mi juicio, su
lectura debe ser cuidadosa: conservar la titularidad no significa que
todo ejercicio de un derecho tenga identicas condiciones o que cualquier
procedimiento pueda iniciarse sin atender a sus reglas.

El iusnaturalismo situa el fundamento en la condicion humana, con
independencia de su reconocimiento positivo. Me resulta convincente
su capacidad para cuestionar una norma injusta: si la autoridad fuera
la unica fuente de valor, seria dificil explicar por que una persona
merece proteccion frente a la propia ley. Su dificultad es justificar
que contenidos se consideran inherentes y con que razones pueden
defenderse ante quienes no comparten el mismo punto de partida.

La perspectiva positivista dirige la atencion hacia el reconocimiento
juridico. Su aportacion es la certeza: permite identificar obligaciones
y distinguir una aspiracion moral de una facultad exigible. No obstante,
considero insuficiente detener el analisis en la existencia de una
norma. Que un derecho este escrito no demuestra que haya medios
adecuados para ejercerlo ni que las autoridades cumplan sus deberes.
La comparacion con el iusnaturalismo es fecunda cuando evita convertir
una diferencia teorica en una eleccion entre dignidad y reglas.

La fundamentacion axiologica entiende los derechos como exigencias
morales previas que orientan la normatividad. El video recoge la critica
de que esta postura puede confundir moral y derecho. Encuentro en esa
objecion una advertencia importante, no una razon para descartar toda
valoracion moral: un valor puede justificar cambiar una ley sin que
esa demanda sea ya una facultad reconocida. Mantener separados ambos
planos permite argumentar con mayor precision.

Valoro la leccion como una introduccion clara a la relacion entre
fundamento, reconocimiento y proteccion de los derechos humanos.
\end{resenabox}

\section{Conclusiones}

Recomiendo el video por distinguir fundamento, reconocimiento y
proteccion de los derechos. Su alcance introductorio exige complementarlo
con Unidad 1 y fuentes mexicanas, sin atribuirle respuestas a casos
que no desarrolla \citep{unedVideo2024,cuadernilloADH2026}.

En Mexico, el articulo 1 exige proteger derechos constitucionales y
convencionales \citep[art.~1]{cpeumADH2026}; la jurisprudencia
1a./J.~37/2016 reconoce la dignidad como derecho fundamental, no mera
declaracion etica \citep{scjnDignidad2016ADH}. P./J.~20/2014 establece
el parametro de regularidad con la salvedad de restricciones
constitucionales expresas \citep{scjnParametro2014ADH}. Estos criterios
impiden confundir una justificacion moral con una respuesta juridica
automatica \citep{ibarraCurso2022ADH}.

Considero que la dignidad permite cuestionar normas excluyentes, pero
requiere obligaciones claras y medios de defensa. Para mi formacion
en Derecho, esa distincion permite exigir proteccion efectiva frente
al abuso, sin confundir legalidad formal con justicia.\footnote{Se utilizo
GitHub Copilot para apoyar la organizacion, el contraste documental y
la redaccion; no sustituye el analisis propio ni la revision personal.}

\clearpage
\begin{thebibliography}{7}
\bibitem[Alvarez Icaza Longoria(2009)]{alvarezIcaza2009ADH}
Alvarez Icaza Longoria, E. (2009). \emph{Los derechos humanos en Mexico}.
Nostra Ediciones.

\bibitem[Ibarra Olguin(2022)]{ibarraCurso2022ADH}
Ibarra Olguin, A. M. (Ed.). (2022). \emph{Curso de derechos humanos}.
Tirant lo Blanch.

\bibitem[Suprema Corte de Justicia de la Nacion, Pleno(2014)]{scjnParametro2014ADH}
Suprema Corte de Justicia de la Nacion, Pleno. (2014).
\emph{Derechos humanos contenidos en la Constitucion y en los tratados
internacionales. Constituyen el parametro de control de regularidad
constitucional, pero cuando en la Constitucion haya una restriccion
expresa al ejercicio de aquellos, se debe estar a lo que establece el
texto constitucional} [Jurisprudencia P./J. 20/2014 (10a.), registro 2006224].
Gaceta del Semanario Judicial de la Federacion, libro 5, tomo I, 202.
\url{https://sjf2.scjn.gob.mx/detalle/tesis/2006224}

\bibitem[Suprema Corte de Justicia de la Nacion, Primera Sala(2016)]{scjnDignidad2016ADH}
Suprema Corte de Justicia de la Nacion, Primera Sala. (2016).
\emph{Dignidad humana. Constituye una norma juridica que consagra un
derecho fundamental a favor de las personas y no una simple declaracion
etica} [Jurisprudencia 1a./J. 37/2016 (10a.), registro 2012363].
Gaceta del Semanario Judicial de la Federacion, libro 33, tomo II, 633.
\url{https://sjf2.scjn.gob.mx/detalle/tesis/2012363}

\bibitem[Camara de Diputados del H. Congreso de la Union(2026)]{cpeumADH2026}
Camara de Diputados del H. Congreso de la Union. (2026).
\emph{Constitucion Politica de los Estados Unidos Mexicanos}
(ultima reforma DOF 7 de octubre de 2026).
\url{https://www.diputados.gob.mx/LeyesBiblio/pdf/CPEUM.pdf}

\bibitem[Lozano Ortega(2026)]{cuadernilloADH2026}
Lozano Ortega, D. A. (2026). \emph{Unidad 1. Bases teoricas basicas
relacionadas con los derechos humanos en Mexico}
[Cuadernillo de contenidos, semana 2]. Universidad Abierta y a Distancia de Mexico.
\url{https://aulavirtual.unadmexico.mx/mod/resource/view.php?id=211995}

\bibitem[Universidad Nacional de Educacion a Distancia(2024)]{unedVideo2024}
Universidad Nacional de Educacion a Distancia. (2024, 23 de octubre).
\emph{Tema 11. Los Derechos Humanos Concepto y Fundamentacion
(1er Curso, 1er Cuatrimestre del GRADO EN DERECHO)} [Video]. Canal UNED.
\url{https://canal.uned.es/video/671b6728e0637d00c70d7b80}
\end{thebibliography}

\end{document}